\input{../tex-inputs/latex-leseansicht-vorspann}

\section[Arthur und Olga Schnitzler an Gustav Schwarzkopf, {[}5?{]}. 9. 1907]{L04374 Arthur und Olga Schnitzler an Gustav Schwarzkopf, {[}5?{]}. 9. 1907}
\nopagebreak\mylabel{L04374v}
\rehead{ }\normalsize\beginnumbering\briefempfaengerindex{Schwarzkopf, Gustav@\textsc{Schwarzkopf, Gustav}!zzzSchnitzler, Olga@\emph{von Olga Schnitzler}!1907-09-051@{{[}5?{]}. 9. 1907}|(be}\briefempfaengerindex{Schwarzkopf, Gustav@\textsc{Schwarzkopf, Gustav}!zzzSchnitzler, Arthur@\emph{von Arthur Schnitzler}!1907-09-051@{{[}5?{]}. 9. 1907}|(be}
\toendnotes[C]{\smallbreak\pagebreak[2]}
\correspDesc{Versand  durch Arthur Schnitzler, Olga Schnitzler am [5?]. 9. 1907 in Meran
\newline{}Übermittlung  am 5. 9. 1907 in Meran
\newline{}Erhalt  durch Gustav Schwarzkopf im Zeitraum [6. 9. 1907 – 10. 9. 1907?] in Wien}\toendnotes[C]{\smallbreak}
\Standort{DLA, A:Schnitzler, HS.1985.1.1897.}
\physDesc{Bildpostkarte, 178 Zeichen
\newline{}Handschrift Arthur Schnitzler: Bleistift, deutsche Kurrent
\newline{}Handschrift Olga Schnitzler: Bleistift, lateinische Kurrent
\newline{}Versand: Stempel: »\nobreak{}\oindex{Meran@\textbf{Meran}|pwk}\textcolor{gray}{M}eran, 5. IX. 7, 7\nobreak{}«.  }\toendnotes[C]{\smallbreak}\pstart{}{\pb}\textsc{Herrn Gustav
                     Schwarzkopf}\pend{}\pstart{}Wien I\oindex{I., Innere Stadt@\textbf{I., Innere Stadt}|pw}\pend{}\pstart{}\textsc{Tiefer Graben 23}\oindex{Wien@\textbf{Wien}!I., Innere Stadt@\textbf{I., Innere Stadt}!Tiefer Graben 23@\textbf{Tiefer Graben 23}|pw}\pend{}{\bigskip}
\pstart
           \centering{}{\pb}\textcolor{gray}{\textbf{Meran\oindex{Meran@\textbf{Meran}|pw} gegen d. Mendel\oindex{Mendelpass@\textbf{Mendelpass}|pw}}}\pend
           \vspace{1em}
\pstart
           \noindent{}{\pb}{[}hs. O. Schnitzler:{]} Herzliche Grüsse,\pend
           \pstart \spacefill\mbox{Olga Schnitzler.}\pend{}\selectlanguage{ngerman}\vspace{1em}
\pstart
           \noindent{}{[}hs. O. Schnitzler:{]} Auf baldgs Wiederſehen, wir denken \label{K_L04374-1v}\edtext{Ende nächſter Woche}{\lemma{\textnormal{\emph{Ende nächster Woche}}}\Cendnote{\textnormal{Die Rückkehr nach Wien\oindex{Wien@\textbf{Wien}|pwk} erfolge nach über zwei Monaten der Abwesenheit am 13. 9. 1907.}}}\label{K_L04374-1} heimzukehren. Es iſt wundervoll hier.\pend
           
\pstart
           Ihr{\\[\baselineskip]}\spacefill\mbox{Arthur}\pend
           \leftskip=0em{}\selectlanguage{ngerman}\endnumbering\briefempfaengerindex{Schwarzkopf, Gustav@\textsc{Schwarzkopf, Gustav}!zzzSchnitzler, Olga@\emph{von Olga Schnitzler}!1907-09-051@{{[}5?{]}. 9. 1907}|)be}\briefempfaengerindex{Schwarzkopf, Gustav@\textsc{Schwarzkopf, Gustav}!zzzSchnitzler, Arthur@\emph{von Arthur Schnitzler}!1907-09-051@{{[}5?{]}. 9. 1907}|)be}\mylabel{L04374h}
\begin{anhang}
\end{anhang}\newcommand{\dateiname}{L04374}\newcommand{\titel}{Arthur und Olga Schnitzler an Gustav Schwarzkopf, [5?]. 9. 1907}\newcommand{\editorInnen}{Herausgegeben von Jahnke, SelmaMüller, Martin Anton}\input{../tex-inputs/latex-leseansicht-abspann}
